\section{Porównanie układów i~rozliczanie zakładów}
\label{sec:showdown-i-rozliczenie}

Porównanie układów (ang. \emph{showdown}) jest fazą końcową każdego rozdania, w~którym gracz wykonał zakład Play.
Środowisko odkrywa wszystkie brakujące karty wspólne, ocenia najlepszą
pięciokartową rękę gracza i~krupiera, a~następnie porównuje otrzymane
wartości. W~przypadku kiedy gracz spasuje, nie dochodzi do porównania układów.

\subsection{Przebieg porównania układów}
\label{subsec:przebieg-showdownu}

Podczas porównania układów gracz i~krupier dysponują siedmioma kartami.
Ewaluator jest odpowiedzialny za wybór najlepszego pięciokartowego układu osobno dla gracza
i~krupiera. Wynik tej operacji jest przechowywany w~obiekcie
\texttt{ShowdownResult}, który zawiera:

\begin{itemize}
    \item najlepszą rękę gracza,
    \item najlepszą rękę krupiera,
    \item wynik porównania z~perspektywy gracza,
    \item informację o~kwalifikacji krupiera.
\end{itemize}

Zbiór wszystkich możliwych wyników porównania układów przedstawiono następująco:

\begin{equation}
    \mathcal{O}_{showdown}
    =
    \{
        \texttt{PLAYER\_WIN},
        \texttt{DEALER\_WIN},
        \texttt{PUSH}
    \}.
\end{equation}

Konkretny rezultat porównania rąk gracza i~krupiera jest wyznaczany przez funkcję \(O\):
dla wartości rąk \(H_p\) oznaczającej rękę gracza oraz \(H_d\) oznaczającej rękę krupiera, wynik porównania układów można zdefiniować jako:

\begin{equation}
    O(H_p,H_d)
    =
    \begin{cases}
        \texttt{PLAYER\_WIN}, & H_p > H_d, \\
        \texttt{DEALER\_WIN}, & H_p < H_d, \\
        \texttt{PUSH},        & H_p = H_d.
    \end{cases}
\end{equation}

Wydawać by się mogło, że w~momencie kiedy krupier się nie kwalifikuje, nie ma sensu porównywać rąk obu stron.
Można jednak rozważyć przykład, w~którym każda ze stron ma jedynie wysokie karty (najniższy możliwy układ),
ale najwyższa karta krupiera jest wyższa od najwyższej karty gracza. W~takim przypadku krupier wygrywa rozdanie, mimo że się nie kwalifikuje.
Dlatego porównanie ręki gracza i~krupiera jest wykonywane nawet w~sytuacji, kiedy krupier się nie kwalifikuje.

\subsection{Kwalifikacja krupiera}
\label{subsec:kwalifikacja-krupiera}

W~Ultimate Texas Hold’em krupier kwalifikuje się do wypłaty zakładu Ante,
jeżeli posiada co najmniej jedną parę, taki warunek można zapisać w~formie równania:

\begin{equation}
    Q(H_d)
    =
    \begin{cases}
        1, & rank(H_d) \geq \texttt{ONE\_PAIR}, \\
        0, & rank(H_d) = \texttt{HIGH\_CARD}.
    \end{cases}
\end{equation}

gdzie \(rank(H_d)\) oznacza rangę układu krupiera.

Samo zjawisko kwalifikacji krupiera nie zmienia wyniku porównania rąk, jednakże wpływa na sposób
rozliczenia zakładu Ante, ponieważ:

\begin{itemize}
    \item jeżeli krupier kwalifikuje się, zakład Ante wygrywa lub przegrywa zgodnie
          z~wynikiem porównania układów,
    \item jeżeli krupier nie kwalifikuje się, zakład Ante jest zwracany niezależnie
          od tego, która strona ma silniejszą rękę.
\end{itemize}

Zakład Play jest rozliczany niezależnie od kwalifikacji krupiera, w~wyniku bezpośredniego porównania rąk.
Z~kolei zakład Blind jest rozliczany według układu gracza (w~przypadku jego
zwycięstwa), przy remisie zwracany, a~przy przegranej gracza jest przegrywany.

Wszystkie wyżej opisane zależności przedstawiono dla ułatwienia w~tabeli~\ref{tab:rozliczenie-showdownu}.

\begin{table}[htbp]
    \centering
    \caption{Rozliczenie zakładów po porównaniu układów}
    \label{tab:rozliczenie-showdownu}
    {\setlength{\tabcolsep}{3pt}
    \begin{tabular}{
        |>{\centering\arraybackslash}m{0.20\textwidth}
        |>{\centering\arraybackslash}m{0.13\textwidth}
        |>{\centering\arraybackslash}m{0.17\textwidth}
        |>{\centering\arraybackslash}m{0.20\textwidth}
        |>{\centering\arraybackslash}m{0.17\textwidth}|
    }
        \hline
        \textbf{Wynik}
        &
        \textbf{Kwalifikacja}
        &
        \textbf{Ante}
        &
        \textbf{Blind}
        &
        \textbf{Play}
        \\
        \hline

        Wygrana gracza
        &
        tak
        &
        wygrana \(1{:}1\)
        &
        według tabeli wypłat
        &
        wygrana \(1{:}1\)
        \\
        \hline

        Wygrana gracza
        &
        nie
        &
        zwrot
        &
        według tabeli wypłat
        &
        wygrana \(1{:}1\)
        \\
        \hline

        Remis
        &
        dowolna
        &
        zwrot
        &
        zwrot
        &
        zwrot
        \\
        \hline

        Wygrana krupiera
        &
        tak
        &
        przegrana
        &
        przegrana
        &
        przegrana
        \\
        \hline

        Wygrana krupiera
        &
        nie
        &
        zwrot
        &
        przegrana
        &
        przegrana
        \\
        \hline
    \end{tabular}}
\end{table}

\subsection{Model rozliczenia zakładów}
\label{subsec:model-rozliczenia}

Zaprojektowano obiekt \texttt{WagerSettlement}, który reprezentuje rozliczenie pojedynczego zakładu.
Składa się on z~\texttt{stake}, czyli początkowej wartości postawionego zakładu, oraz \texttt{net\_profit}, czyli zysku albo straty netto.

Dla pojedynczego zakładu
o~stawce \(s\) i~zysku netto \(n\) całkowity zwrot brutto wynosi
\(g = s + n\), jeśli \(g = s\) oznacza to zwrot stawki, w~przypadku \(g = 0\) jej pełną utratę, a~stawka \(s = 0\) zakład niepostawiony.

Pełne rozliczenie rozdania reprezentuje obiekt \texttt{Settlement},
zawierający osobne wyniki dla Ante, Blind oraz Play.
Łączny wynik netto \(S_{net}\) jest sumą zysków netto wszystkich trzech zakładów, można go zapisać w~postaci równania:

\begin{equation}
    S_{net}
    =
    n_{Ante}
    +
    n_{Blind}
    +
    n_{Play}.
\end{equation}

Do reprezentacji wartości finansowych wykorzystano typ \texttt{Fraction}, ponieważ ten typ pozwala na
przechowywanie wypłaty takiej jak \(3/2\) (bez stosowania przybliżeń zmiennoprzecinkowych), dzięki czemu
unika się błędów zaokrągleń.


\subsection{Rozliczenie zakładu Blind}
\label{subsec:rozliczenie-blind}

Zakład Blind ma stałą stawkę równą jednej jednostce Ante. W~sytuacji kiedy wygrywa gracz,
zysk z~racji tego zakładu jest uzależniony od tego, jak wysoki układ uda się graczowi osiągnąć.
Dla układów słabszych od strita Blind jest zwracany. Dla strita lub
silniejszego układu stosowana jest tabela wypłat przedstawiona w~
tabeli~\ref{tab:wyplaty-blind-implementacja}.

\begin{table}[htbp]
    \centering
    \caption{Tabela wypłat zakładu Blind}
    \label{tab:wyplaty-blind-implementacja}
    \begin{tabular}{
        |m{0.30\textwidth}
        |>{\centering\arraybackslash}m{0.15\textwidth}
        |>{\centering\arraybackslash}m{0.15\textwidth}|
    }
        \hline
        \textbf{Układ gracza}
        &
        \textbf{Wypłata}
        &
        \textbf{Zysk netto}
        \\
        \hline

        Wysoka karta
        &
        zwrot
        &
        \(0\)
        \\
        \hline

        Jedna para
        &
        zwrot
        &
        \(0\)
        \\
        \hline

        Dwie pary
        &
        zwrot
        &
        \(0\)
        \\
        \hline

        Trójka
        &
        zwrot
        &
        \(0\)
        \\
        \hline

        Strit
        &
        \(1{:}1\)
        &
        \(1\)
        \\
        \hline

        Kolor
        &
        \(3{:}2\)
        &
        \(\frac{3}{2}\)
        \\
        \hline

        Full
        &
        \(3{:}1\)
        &
        \(3\)
        \\
        \hline

        Kareta
        &
        \(10{:}1\)
        &
        \(10\)
        \\
        \hline

        Poker
        &
        \(50{:}1\)
        &
        \(50\)
        \\
        \hline

        Poker królewski
        &
        \(500{:}1\)
        &
        \(500\)
        \\
        \hline
    \end{tabular}
\end{table}

Zakład Blind jest stosowany wyłącznie wtedy, gdy gracz wygrał porównanie układów.
Przy remisie zakład jest zwracany, natomiast przy wygranej krupiera
pełna stawka Blind zostaje utracona.

\subsection{Rozliczenie Ante, Play i~spasowania}
\label{subsec:rozliczenie-ante-play}

Zakład Ante przy remisie jest zwracany.
Jeżeli krupier kwalifikuje się, Ante wygrywa albo przegrywa zgodnie
z~wynikiem porównania układów. W~przypadku braku kwalifikacji zakład jest zwracany.

Zakład Play jest wypłacany \(1{:}1\), niezależnie od jego mnożnika.
Dla zobrazowania można posłużyć się przykładem: gracz postawił
wygrany zakład Play \(4\times\), co daje zysk netto równy czterem
jednostkom. Przy porażce tracona jest pełna wartość zakładu, natomiast
przy remisie cała stawka zostaje zwrócona.

Jeżeli gracz spasuje na riverze, nie dochodzi do porównania układów. Ante i~Blind
zostają przegrane, natomiast Play nie został postawiony.

Rozliczenie spasowania ma zatem postać

\begin{equation}
    S_{fold}
    =
    \left(
        n_{Ante}=-1,
        n_{Blind}=-1,
        n_{Play}=0
    \right).
\end{equation}

Wtedy łączna stawka wynosi dwie jednostki, a~wynik netto jest równy \(-2\).

W~celu rozróżnienia zakładu niepostawionego od zakładu postawionego,
ale zwróconego, w~rozliczeniu Play zapisywana jest stawka równa zero.

\section{Funkcja nagrody i~wynik epizodu}
\label{sec:funkcja-nagrody}

Aby móc sprawnie zmieniać warianty funkcji nagrody wykorzystywanej przez algorytmy uczenia przez wzmacnianie,
oddzielono mechanizm rozliczania zakładów od konstrukcji sygnału nagrody.
Silnik domenowy
wyznacza dokładny wynik finansowy rozdania zgodny z~zasadami
Ultimate Texas Hold’em. Natomiast
konstruowanie sygnału nagrody na jego podstawie należy do warstwy
integrującej silnik z~agentem. Dzięki temu możliwe jest proste porównywanie różnych funkcji nagrody, bez
ingerencji w~zasady gry.
